\section{Separate input traversal from retained output}
Two complete private prototypes use different storage strategies. The candidate version
keeps decoded payloads in a private record until validation finishes. The checked-slice
version retains input locations, validates the whole record and then revisits the sample
heads during publication. Its generic walker uses explicit frames; the candidate
version's walker uses recursive calls. Both produce identical required-field encodings.

The table records completed MIPS instructions inside \texttt{telemetry\_decode}, from
entry until its return to the caller. The inputs are the helper's trace, maximum-output
record and exact-input-budget record. Results are identical at ordinary and above-4-GiB
addresses. These figures are observed work, not elapsed host time or grading targets.

\tagpdfsetup{table/header-rows={1}}
\begin{center}
\begin{tabular}{lrrrr}
Input & Input bytes & Output bytes & Candidate & Checked slices \\
\hline
Worked trace & 52 & 37 & 2186 & 2620 \\
Maximum output & 241 & 241 & 7199 & 11344 \\
Exact input budget & 1024 & 7 & 7712 & 5454 \\
\end{tabular}
\end{center}

The maximum-output record has sixty-four readings. Revisiting those heads adds work
during publication. The input-budget record has empty required payloads and a large
optional array. Its output is only seven bytes, but the decoder must still walk and
validate the optional children. Dropping metadata does not authorize jumping past an
unvalidated item. The explicit-frame prototype performs less work on this particular
container-heavy input, while it performs more work on the other two.

These prototypes also differ in their head readers and size calculations. The table
therefore does not isolate the cost of recursion or a second pass. It establishes that
output size alone does not predict decoding work, and that one input cannot establish
which complete method does less work in general. Both methods must meet the same field,
range, depth and publication contracts before their work counts are compared.

The table uses the earlier complete prototype sources. The cumulative construction
adds caller routes, checkpoint instructions and unfinished-child guards. Those changes
make its executed instruction counts different, even for the same field values; do
not substitute them for the measurement boundaries recorded here.
